\documentclass[11pt]{article}

\usepackage[a4paper,margin=1in]{geometry}
\usepackage{amsmath,amssymb,amsthm}
\usepackage{enumitem}
\usepackage{microtype}
\usepackage{parskip}

% -------------------------------------------------------------------
% Environments
% -------------------------------------------------------------------
\theoremstyle{definition}
\newtheorem*{definition}{Definition}
\newtheorem*{notation}{Notation}
\newtheorem*{example}{Example}

% -------------------------------------------------------------------
% Notation
% -------------------------------------------------------------------
\newcommand{\DB}{D_B}
\newcommand{\dB}{d_B}

% A problem cell: \cell{<number>}{<title>}{<points>}
\newcommand{\cell}[3]{%
  \section*{Cell #1: #2\hfill{\normalsize\mdseries #3}}}

\title{Bulgarian Solitaire\\[0.4em]\large Problem Description}
\author{}
\date{}

\begin{document}
\maketitle

% ===================================================================
\section*{Overview}
% ===================================================================

Repeatedly take one card from every pile and form a new pile.
How many moves can it take before the piles start to cycle?

\begin{definition}[Partition]
A \emph{partition} of a positive integer $n$ is a weakly decreasing
sequence
\[
  \lambda = (\lambda_1, \ldots, \lambda_s)
\]
of positive integers with
\[
  \lambda_1 + \cdots + \lambda_s = n.
\]
Think of it as a division of $n$ cards into $s$ piles.
\end{definition}

\begin{definition}[Shift]
The \emph{shift} $B(\lambda)$ is the partition of $n$ obtained as
follows: remove one card from every pile, discard the piles that
become empty, and add one new pile consisting of the $s$ removed
cards. Formally, $B(\lambda)$ is the partition whose parts are the
positive numbers among
\[
  \lambda_1 - 1, \; \ldots, \; \lambda_s - 1
\]
together with one extra part equal to $s$.
\end{definition}

Iterating $B$ from any starting partition must eventually repeat,
since there are finitely many partitions of $n$.

\begin{definition}[Cyclic partitions and time to cycle]
Call $\lambda$ \emph{cyclic} if
\[
  B^i(\lambda) = \lambda \qquad \text{for some } i \geq 1,
\]
and let
\begin{align*}
  \dB(\lambda) &= \min\bigl\{\, i \geq 0 : B^i(\lambda) \text{ is cyclic} \,\bigr\}, \\
  \DB(n)       &= \max\bigl\{\, \dB(\lambda) : \lambda \text{ is a partition of } n \,\bigr\}.
\end{align*}
Thus $\DB(n)$ is the largest number of shifts that can be needed
before the process starts to cycle.
\end{definition}

\begin{notation}[Triangular numbers and rank]
Write
\[
  T_k = \frac{k(k+1)}{2},
  \qquad
  \delta_k = (k, k-1, \ldots, 2, 1),
\]
the partition of $T_k$ into distinct parts. Every $n \geq 1$
satisfies
\[
  T_{k-1} < n \leq T_k
\]
for exactly one $k$, which we call the \emph{rank} of $n$.
\end{notation}

\begin{example}
\begin{align*}
  B\bigl((2,1,1,1,1)\bigr) &= (5,1), \\
  B\bigl((5,1)\bigr)       &= (4,2), \\
  B\bigl((4,2)\bigr)       &= (3,2,1), \\
  B\bigl((3,2,1)\bigr)     &= (3,2,1).
\end{align*}
Here $\dB\bigl((2,1,1,1,1)\bigr) = 3$.
\end{example}

% ===================================================================
\section*{What to Hand In}
% ===================================================================

For each cell you attempt, hand in a written proof.

\begin{itemize}[leftmargin=*]
  \item \textbf{Computation.}
    You may use a computer to explore. A proof may rely on a
    computation only if you include the code, it runs in under
    10 minutes on a laptop, and the computation is exhaustive over a
    finite set that your argument has reduced the problem to.
  \item \textbf{Exact values.}
    Where a cell asks you to determine a quantity, give the exact
    value as a formula in $k$ and prove both bounds; state any
    partitions you use explicitly as functions of $k$.
  \item \textbf{Status.}
    Say clearly which cells you consider solved and which are
    partial.
  \item \textbf{Citations.}
    Citing a published result for the statement you are asked to
    prove does not count as a solution; a proof written out in full
    does, whatever its source.
\end{itemize}

% ===================================================================
\cell{1}{Cyclic Partitions and Cycles}{}
% ===================================================================

\emph{The statement of this cell is not visible in the provided
screenshots.}

% ===================================================================
\cell{2}{$\DB$ at Triangular $n$}{2 points}
% ===================================================================

Next, how long the process can take to reach a cycle, starting with
the triangular numbers. Determine
\[
  \DB(T_k)
\]
for every $k$, with proof of both bounds.

% ===================================================================
\cell{3}{A General Upper Bound}{3 points}
% ===================================================================

Now the numbers strictly between two consecutive triangular numbers.

\begin{enumerate}[label=(\alph*), leftmargin=*]
  \item Prove that for every $k \geq 4$ and every non-triangular $n$
    with $T_{k-1} < n < T_k$,
    \[
      \DB(n) \leq k^2 - 2k - 1.
    \]
  \item Determine $\DB(T_k - 1)$ exactly.
  \item Determine also, for that $n$, which partitions attain the
    maximum.
\end{enumerate}

% ===================================================================
\cell{4}{One Above a Triangular Number}{5 points}
% ===================================================================

The first family just above a triangular number:
\[
  n = T_{k-1} + 1,
\]
that is,
\[
  n = 11, 16, 22, 29, \ldots
  \qquad \text{for } k = 5, 6, 7, 8, \ldots.
\]
Determine
\[
  \DB(T_{k-1} + 1)
\]
for every $k \geq 5$, with proof of both bounds.

% ===================================================================
\cell{5}{Two Above a Triangular Number}{8 points}
% ===================================================================

The next family:
\[
  n = T_{k-1} + 2,
\]
that is,
\[
  n = 3, 5, 8, 12, 17, 23, \ldots
  \qquad \text{for } k = 2, 3, 4, 5, 6, 7, \ldots.
\]
Determine
\[
  \DB(T_{k-1} + 2)
\]
for every $k$, with proof of both bounds.

\begin{itemize}[leftmargin=*]
  \item State exactly for which $k$ your formula holds, and give the
    remaining values separately.
  \item Your upper bound must be a single argument valid for all $k$
    in that range, not a separate treatment of each $k$.
  \item You must give the extremal partitions explicitly as a
    function of $k$.
  \item Say also where the straightforward extension of the Cell~3
    argument stops: give the bound it does yield, show it is strictly
    weaker than the truth, and identify precisely what your proof
    supplies in its place.
\end{itemize}

% ===================================================================
\cell{6}{$\DB(n)$ for Every $n$}{13 points}
% ===================================================================

\textbf{Open question.}

Finally, the whole function
\[
  n \longmapsto \DB(n).
\]
Determine $\DB(n)$ for every $n$.

\end{document}
